% !TEX TS-program = xelatex
\documentclass[12pt,a4paper]{article}
\usepackage[utf8]{inputenc}
\usepackage{geometry}
\geometry{margin=25mm}
\usepackage{amsmath, amssymb}
\usepackage{booktabs}
\usepackage{graphicx}
\usepackage{hyperref}
\usepackage{siunitx}
\usepackage{enumitem}
\usepackage{xcolor}
\usepackage{csquotes}

\hypersetup{
  colorlinks=true,
  linkcolor=black,
  urlcolor=blue,
  citecolor=black
}

\title{Analyzing Bias in Word Embeddings: A WEAT-Based Approach}
\date{\today}

\begin{document}
\maketitle

\begin{abstract}
This study implements the Word Embedding Association Test (WEAT) to analyze 
social biases present in pre-trained Word2Vec embeddings. 
In the experiment with flowers versus insects and positive versus negative attributes, 
a strong bias was observed with effect size $d > 0.8, p < 0.01$. 
In the gender--career/family test, male-related words were strongly associated with career terms, 
while female-related words were associated with family terms. 
In the gender--science/arts test, only a weak tendency was found. 
The interpretation of effect sizes follows Cohen (1988), where 
$d \approx 0.2$ is considered a small effect, 
$d \approx 0.5$ a medium effect, 
and $d \geq 0.8$ a large effect. 
These results suggest that embeddings inherently reflect social stereotypes, 
which can influence downstream applications.
\end{abstract}

\section{Introduction (Background \& Related Works)}
Pre-trained embeddings learn statistical patterns from large corpora, 
which can also encode social stereotypes. 
\citet{caliskan2017} demonstrated through WEAT that word embeddings reflect gender and racial biases. 
Subsequent works, such as \citet{bolukbasi2016} and \citet{may2019}, 
extended this line of research to sentence-level evaluations and proposed debiasing methods. 
In this study, we implement WEAT on Word2Vec embeddings to quantify the degree of bias.

\section{Method}
\subsection{Model}
We used the GoogleNews Word2Vec model (300 dimensions, binary format, limited to the top 500k words). 
Additionally, we trained a Word2Vec model on Korean tokenized data 
(100 dimensions, window size = 5, min\_count = 3).

\subsection{Target and Attribute sets}
\begin{itemize}
  \item Flowers (rose, tulip, lily, daisy) vs Insects (spider, mosquito, fly, grasshopper)
  \item Positive (love, happiness, laughter) vs Negative (disaster, pain, hatred)
  \item Gender--career/family, Gender--science/arts (WEAT T7, T8)
\end{itemize}

\subsection{Evaluation}
Association scores $s(w,A,B)$ were computed using cosine similarity. 
Effect sizes (Cohen’s $d$) and permutation test $p$-values (10,000 iterations) were calculated. 
Interpretation of effect sizes follows Cohen (1988): 
$d \approx 0.2$ is a small effect, 
$d \approx 0.5$ is a medium effect, 
and $d \geq 0.8$ is a large effect.

\section{Results}
\subsection{Flowers vs Insects -- Positive vs Negative}
Flower words were more closely associated with positive attributes, 
while insect words were associated with negative attributes. 
A strong bias was observed ($d > 0.8$, $p < 0.01$).

\subsection{Gender--Career/Family (WEAT T7)}
Male-related words were more strongly associated with career, 
while female-related words were closer to family terms.  
GoogleNews Word2Vec result: $d \approx 0.8$, $p < 0.01$.

\subsection{Gender--Science/Arts (WEAT T8)}
Male-related words aligned more closely with science, 
while female-related words were closer to arts.  
The effect size was $d \approx 0.4$, $p \approx 0.07$, 
showing only a non-significant trend.

\subsection{Visualization}
When projecting target and attribute words using t-SNE, 
a clear group separation was observed.

\section{Discussion \& Conclusion}
This study implemented WEAT to measure and analyze social biases in Word2Vec embeddings. 
Key findings are as follows:
\begin{enumerate}[leftmargin=2em]
  \item Flowers vs Insects: strong bias observed ($d > 0.8$, $p < 0.01$)
  \item Gender--Career/Family: significant bias ($d \approx 0.8$, $p < 0.01$)
  \item Gender--Science/Arts: non-significant trend ($d \approx 0.4$, $p \approx 0.07$)
\end{enumerate}

These findings indicate that embeddings reflect and may amplify social stereotypes in downstream systems. 
Future work will extend to Korean embeddings, evaluate debiasing techniques, 
and develop specialized benchmarks for bias assessment.

\begin{thebibliography}{9}
\bibitem{caliskan2017}
A.~Caliskan, J.~J. Bryson, and A.~Narayanan.
\newblock Semantics derived automatically from language corpora contain human-like biases.
\newblock \emph{Science}, 356(6334):183--186, 2017.

\bibitem{bolukbasi2016}
T.~Bolukbasi, K.-W. Chang, J.~Zou, V.~Saligrama, and A.~Kalai.
\newblock Man is to Computer Programmer as Woman is to Homemaker? Debiasing Word Embeddings.
\newblock In \emph{Proceedings of NeurIPS}, 2016.

\bibitem{cohen1988}
J.~Cohen.
\newblock Statistical power analysis for the behavioral sciences.
\newblock \emph{Lawrence Erlbaum Associates}, 2nd edition, 1988.
\end{thebibliography}

\end{document}
